\chapter{La nozione di verità e di conseguenza logica}

\textit{Lucidi: \texttt{lucidi/2.Truth}. Approfondimento: Van Dalen,
\emph{Logic and Structure}, §1.1 (Propositions and Connectives) e §1.2
(Semantics).}

% =====================================================================
\section{Che cos'è la verità}
% =====================================================================

Abbiamo visto che la logica è lo studio sistematico della forma delle
\textbf{inferenze valide} e delle leggi più generali della \textbf{verità}.
Ma che cos'è la verità?

\dfn{Verità (senso comune)}{
  La verità è l'essere \textbf{in accordo con i fatti o con la realtà}
  (oppure la fedeltà a un originale o a uno standard).
}

Nelle scienze sperimentali (fisica, chimica, biologia) la verità è associata
al \textbf{mondo sensibile}:
\begin{itemize}
  \item un \textbf{esperimento ripetibile} definisce una verità;
  \item se il mondo sensibile cambia, cambia anche la verità.
\end{itemize}

\qs{Che cosa succede quando non c'è un mondo sensibile, ad esempio in
matematica?}{}

\subsection{La verità in matematica}

Le nostre ``verità probabili'' sono, ad esempio:
\begin{itemize}
  \item la somma degli angoli interni di un triangolo è $180^\circ$;
  \item in modo equivalente, il \textbf{V postulato} di Euclide: dati una retta
        e un punto fuori di essa, esiste \textbf{una e una sola} retta per il
        punto parallela alla retta data.
\end{itemize}
Però\dots

\ex{Geometria sferica (Riemann, 1854)}{
  Reinterpretiamo i termini primitivi della geometria:
  \begin{itemize}
    \item \textbf{punto} = punto sulla superficie di una sfera;
    \item \textbf{retta} = cerchio massimo, cioè intersezione della sfera con
          un piano che passa per il centro.
  \end{itemize}
  Allora:
  \begin{itemize}
    \item dati due punti esiste una sola retta che li contiene;
    \item due rette hanno sempre almeno due punti in comune (quindi
          \emph{non esistono parallele});
    \item la somma degli angoli interni di un triangolo è \textbf{maggiore di
          $180^\circ$}.
  \end{itemize}
}

\nt{
  \textbf{Precisazione (approfondimento).} ``Dati due punti esiste una sola
  retta'' vale per punti \emph{non antipodali}: per due punti opposti (es.\ i
  poli) passano infiniti cerchi massimi. Per questo nella \emph{geometria
  ellittica} si identificano i punti antipodali. Il triangolo con un vertice
  al polo nord e due sull'equatore, distanti un quarto di giro, ha tre angoli
  retti: somma $270^\circ$.
}

\subsection{Conseguenze}

\mprop{Il metodo assiomatico descrive modelli, non il mondo}{
  \begin{itemize}
    \item Il metodo assiomatico \textbf{non è una descrizione del mondo
          sensibile}.
    \item La geometria (la matematica) non sembra fornire giudizi sintetici a
          priori (nel senso di Kant).
    \item Fornisce invece la descrizione di un \textbf{insieme di mondi
          (modelli)} in cui:
      \begin{enumerate}
        \item si possono interpretare le entità (punti, rette, \dots) come si
              vuole, \emph{purché gli assiomi siano rispettati};
        \item mondi diversi hanno proprietà diverse;
        \item solo alcune di queste proprietà sono dimostrabili dagli assiomi
              (\textbf{incompletezza}).
      \end{enumerate}
  \end{itemize}
}

\nt{
  Il punto 3 si legge così: il V postulato è vero nel piano euclideo e falso
  sulla sfera, e \emph{entrambi} soddisfano gli altri assiomi. Quindi il V
  postulato \textbf{non è dimostrabile} dagli altri assiomi (altrimenti sarebbe
  vero in tutti i loro modelli), e neanche la sua negazione. È un enunciato
  \emph{indipendente}. Questo è esattamente il metodo che useremo in logica:
  per mostrare che $F$ \emph{non} segue da $\Gamma$ basta esibire \emph{un}
  mondo in cui $\Gamma$ è vero e $F$ è falso (un \textbf{controesempio}).
}

Per parlarne con precisione introduciamo un po' di lessico.

\dfn{Teoria, modello, consistenza}{
  \begin{itemize}
    \item Una \textbf{teoria} è un insieme di enunciati (assiomi).
    \item Un \textbf{modello} di una teoria è un mondo, con un'interpretazione
          dei simboli, in cui tutti gli assiomi sono veri.
    \item Una teoria è \textbf{consistente} se da essa non si può derivare
          l'assurdo.
  \end{itemize}
}

\mprop{Impatto sulla visione della matematica}{
  \begin{itemize}
    \item Una teoria non si accetta o scarta perché è soddisfatta (banalmente)
          dal mondo sensibile, ma perché \textbf{ammette dei modelli}.
    \item Una teoria che ammette modelli \textbf{non è mai inconsistente}: da
          essa non si può derivare l'assurdo.
    \item La nozione assoluta di verità (data dal mondo sensibile) viene
          sostituita dalla nozione di \textbf{conseguenza logica}.
  \end{itemize}
}

\nt{
  Perché ``ha un modello $\Rightarrow$ consistente''? Le regole di deduzione
  sono \textbf{corrette} (\emph{soundness}): trasformano enunciati veri in un
  mondo in enunciati veri nello stesso mondo. In un modello tutti gli assiomi
  sono veri, quindi è vero tutto ciò che si deriva; ma l'assurdo $\bot$ non è
  vero in nessun mondo, quindi non si deriva. Il viceversa (consistente
  $\Rightarrow$ ha un modello) è il \textbf{teorema di completezza}, che
  vedremo (Van Dalen §1.5).
}

\subsection{Riassunto}

\begin{itemize}
  \item \textbf{Verità del mondo sensibile}: cambia quando cambia il mondo
        (mutamenti \emph{ristretti}: il mondo cambia secondo le sue leggi).
  \item \textbf{Verità della matematica}: cambia adottando nuovi mondi
        (mutamenti \emph{non ristretti}: qualsiasi mondo che rispetta gli
        assiomi).
  \item Inoltre la verità cambia se si \textbf{interpretano} in modo diverso
        le parole di un enunciato.
\end{itemize}

\qs{Che cosa resta della nozione di verità se il mondo cambia arbitrariamente?}{
  Resta ciò che \emph{non dipende} dal mondo: la conseguenza logica.
}

% =====================================================================
\section{Conseguenza ed equivalenza logica}
% =====================================================================

\dfn{Conseguenza logica (lucidi)}{
  Sia $\Gamma = \{F_1, \dots, F_n\}$ un insieme di enunciati (\textbf{ipotesi}
  o \textbf{premesse}) e $F$ un enunciato (\textbf{conclusione}).
  $F$ è \textbf{conseguenza logica} di $\Gamma$, e si scrive
  \[ \Gamma \models F, \]
  se, per tutte le interpretazioni delle formule nei vari mondi possibili,
  vale sempre: \emph{se tutte le formule di $\Gamma$ sono vere (nel mondo in
  esame) allora anche $F$ è vera}.
}

\nt{
  \textbf{Come leggerla.}
  \begin{itemize}
    \item Non si chiede che $F$ sia vera: si chiede che sia vera \emph{in ogni
          mondo in cui le premesse sono vere}.
    \item I mondi in cui qualche premessa è falsa non contano.
    \item $\Gamma \not\models F$ significa che esiste \emph{almeno un} mondo
          (un controesempio) in cui tutte le premesse sono vere e $F$ è falsa.
    \item Con $\Gamma = \emptyset$: $\models F$ significa che $F$ è vera in
          tutti i mondi ($F$ è \textbf{valida}, o \textbf{tautologia} nel caso
          proposizionale).
  \end{itemize}
}

\dfn{Equivalenza logica}{
  $F$ è \textbf{logicamente equivalente} a $G$, e si scrive $F \equiv G$, sse
  $F$ è conseguenza logica di $G$ e $G$ è conseguenza logica di $F$:
  \[ F \equiv G \iff F \models G \ \text{e}\ G \models F. \]
}

\mlenma{Caratterizzazione dell'equivalenza}{
  $F \equiv G$ sse in ogni mondo (con ogni interpretazione) $F$ e $G$ hanno lo
  \textbf{stesso valore di verità}.
}

\pf{Dimostrazione}{
  ($\Rightarrow$) Fissato un mondo: se $F$ è vera, da $F \models G$ anche $G$
  è vera; se $G$ è vera, da $G \models F$ anche $F$ è vera. Quindi $F$ è vera
  sse $G$ è vera: stesso valore.

  ($\Leftarrow$) Se in ogni mondo $F$ e $G$ hanno lo stesso valore, in ogni
  mondo in cui $F$ è vera lo è anche $G$ ($F \models G$) e viceversa
  ($G \models F$).
}

\thm{L'equivalenza logica è una relazione di equivalenza}{
  (Esercizio dei lucidi.) Per tutte le formule $F, G, H$:
  \begin{enumerate}
    \item \textbf{riflessività}: $F \equiv F$;
    \item \textbf{simmetria}: $F \equiv G \Rightarrow G \equiv F$;
    \item \textbf{transitività}: $F \equiv G$ e $G \equiv H \Rightarrow F \equiv H$.
  \end{enumerate}
}

\pf{Dimostrazione}{
  Basta mostrare che $\models$ tra singole formule è riflessiva e transitiva;
  la simmetria di $\equiv$ viene gratis dalla definizione.
  \begin{enumerate}
    \item $F \models F$: in ogni mondo in cui $F$ è vera, $F$ è vera.
          Quindi $F \models F$ e $F \models F$, cioè $F \equiv F$.
    \item $F \equiv G$ significa ``$F \models G$ e $G \models F$'', che è la
          stessa cosa di ``$G \models F$ e $F \models G$'', cioè $G \equiv F$.
    \item Transitività di $\models$: se $F \models G$ e $G \models H$, prendiamo
          un mondo in cui $F$ è vera; allora $G$ è vera (da $F \models G$),
          quindi $H$ è vera (da $G \models H$). Dunque $F \models H$.
          Allo stesso modo da $H \models G$ e $G \models F$ si ha $H \models F$.
          Quindi $F \equiv H$.
  \end{enumerate}
  (In alternativa: per il lemma, $F \equiv G$ sse ``stesso valore in ogni
  mondo'', e ``avere lo stesso valore'' è chiaramente riflessiva, simmetrica e
  transitiva.)
}

\subsection{Conclusioni sulla verità}

\begin{itemize}
  \item La verità è sempre definita \textbf{rispetto a un mondo} e a una data
        \textbf{interpretazione} dei simboli in quel mondo.
  \item La verità di un singolo enunciato \textbf{non è interessante}: cambia
        quando cambia il mondo.
  \item Ciò che interessa sono la \textbf{conseguenza logica} e
        l'\textbf{equivalenza logica}: \textbf{non cambiano} al variare del
        mondo.
\end{itemize}

% =====================================================================
\section{Logica proposizionale}
% =====================================================================

La logica proposizionale è la logica più semplice; illustra le idee di base
usando le \textbf{proposizioni}.

\dfn{Proposizione atomica e valore di verità}{
  Una \textbf{proposizione} è un enunciato \emph{dichiarativo}, a cui si può
  (in linea di principio) assegnare un valore vero o falso. Esempi:
  \begin{itemize}
    \item $P_1$: la neve è bianca;
    \item $P_2$: oggi piove;
    \item $P_3$: questo corso è noioso.
  \end{itemize}
  Ogni $P_i$ che non ha parti che siano a loro volta proposizioni è un
  \textbf{atomo} (o formula atomica).
  Ogni $P_i$ è \textbf{vero o falso, ma mai entrambi} (logica a due valori).
  Il valore vero o falso assegnato a una proposizione è il suo
  \textbf{valore di verità}.
}

\nt{
  (Van Dalen §1.1) Non serve sapere \emph{di fatto} quale sia il valore (la
  congettura di Goldbach è una proposizione anche se non sappiamo se è vera);
  basta che un valore ci sia. Domande e ordini (``chiudi la porta!'') non sono
  proposizioni: non hanno valore di verità.
}

\subsection{I connettivi}

Il linguaggio naturale ha molti connettivi: \emph{and, or, not, but, since,
if\dots then}, \dots\ Perché la logica ne usa pochi?

\ex{``But'' vs.\ ``and''}{
  \begin{itemize}
    \item ``Max è marxista, \emph{ma} non è privo di umorismo.''
    \item ``$\pi$ è irrazionale, \emph{ma} non è algebrico.''
  \end{itemize}
  Nel primo ``ma'' porta un significato emotivo (sorpresa). Nel secondo la
  sorpresa è meno immediata (bisognerebbe sapere che ``quasi tutti'' gli
  irrazionali sono algebrici) e si può usare ``e'' senza cambiare significato.
}

Per \textbf{economia} scegliamo solo connettivi utili per esprimere le
dimostrazioni, come ``e'' e ``o''. Ma anche questi sono ambigui.

\ex{Problema 1: ``and'' e il tempo}{
  \begin{itemize}
    \item John è ripartito \emph{e} ha investito un pedone.
    \item John ha investito un pedone \emph{e} è ripartito.
  \end{itemize}
  In italiano ``e'' può avere una funzione di \textbf{ordinamento temporale}
  (le due frasi dicono cose diverse). In matematica e in logica no:
  $A \wedge B$ significa solo che valgono entrambe, e $A \wedge B$ e
  $B \wedge A$ dicono la stessa cosa.
}

\ex{Problema 2: ``if\dots then'' e ``or''}{
  \begin{enumerate}
    \item Se apro la finestra allora avremo aria fresca.
    \item Se apro la finestra allora $1 + 3 = 4$.
    \item Se $1 + 2 = 4$ allora avremo aria fresca.
    \item John sta lavorando o è a casa.
    \item Euclide era un greco o un matematico.
  \end{enumerate}
  \begin{itemize}
    \item In 1 c'è una \textbf{relazione} tra premessa e conclusione.
    \item In 2 e 3 non c'è alcuna relazione, eppure in matematica sono
          perfettamente accettabili: è l'\textbf{implicazione materiale}.
    \item Tendiamo ad accettare 4 e a rifiutare 5: nel parlato ``o'' suona
          esclusivo, e non si usa ``o'' quando si potrebbe dire ``e''.
          In logica ``o'' è \textbf{inclusivo}: 5 è vera.
  \end{itemize}
}

\mprop{Conclusione}{
  Serve un \textbf{linguaggio artificiale} con un significato \textbf{preciso}
  dei connettivi.
}

\nt{
  \textbf{Perché l'implicazione materiale} (Van Dalen §1.2). ``Se Smith è
  dentro allora Jones è dentro'' è sicuramente \emph{falsa} solo quando Smith
  è dentro e Jones è fuori. Quando Smith è fuori, l'informazione non è falsa,
  è solo ``poco utile''; avendo deciso che ogni enunciato è vero o falso, la
  si dichiara vera. È una \emph{scelta}, motivata dall'eleganza del sistema
  che ne risulta. Da ricordare: \textbf{c'è un solo caso in cui
  un'implicazione è falsa} (premessa vera, conclusione falsa). Quindi 2 è vera
  (conclusione vera) e 3 è vera (premessa falsa).
}

% =====================================================================
\section{Approfondimento: sintassi e semantica (Van Dalen §1.1--1.2)}
% =====================================================================

Formalizziamo ora quanto detto. Si procede in due passi: (1) si definisce il
\textbf{linguaggio} (sintassi); (2) si definisce quando una formula è
\textbf{vera} (semantica). La nozione informale di ``mondo con
un'interpretazione'' diventa, nel caso proposizionale, una
\textbf{valutazione}.

\subsection{Sintassi}

\dfn{Alfabeto della logica proposizionale}{
  \begin{enumerate}[label=(\roman*)]
    \item simboli proposizionali: $p_0, p_1, p_2, \dots$;
    \item connettivi: $\wedge$ (and), $\vee$ (or), $\to$ (implicazione),
          $\neg$ (not), $\leftrightarrow$ (sse), $\bot$ (falsum, l'assurdo);
    \item simboli ausiliari: $($ \ $)$.
  \end{enumerate}
  I simboli proposizionali e $\bot$ sono gli \textbf{atomi}.
}

\dfn{Proposizioni (formule)}{
  L'insieme $\PROP$ è il \textbf{più piccolo} insieme $X$ tale che:
  \begin{enumerate}[label=(\roman*)]
    \item $p_i \in X$ per ogni $i$, e $\bot \in X$;
    \item se $\varphi, \psi \in X$ allora $(\varphi \wedge \psi),
          (\varphi \vee \psi), (\varphi \to \psi), (\varphi \leftrightarrow
          \psi) \in X$;
    \item se $\varphi \in X$ allora $(\neg\varphi) \in X$.
  \end{enumerate}
}

\nt{
  ``Il più piccolo'' significa: in $\PROP$ ci sono \emph{solo} le stringhe
  ottenibili applicando un numero finito di volte le regole (i)--(iii).
  Le lettere greche $\varphi, \psi$ \emph{non} sono formule: sono
  \textbf{metavariabili}, cioè simboli del metalinguaggio che stanno per
  formule qualsiasi. È la separazione linguaggio/metalinguaggio che evita i
  paradossi del capitolo precedente.
}

\thm{Principio di induzione sulle formule}{
  Sia $A$ una proprietà. Se
  \begin{enumerate}[label=(\roman*)]
    \item $A(\varphi)$ vale per ogni atomo $\varphi$,
    \item $A(\varphi)$ e $A(\psi)$ implicano $A(\varphi \square \psi)$ per
          $\square \in \{\wedge, \vee, \to, \leftrightarrow\}$,
    \item $A(\varphi)$ implica $A(\neg\varphi)$,
  \end{enumerate}
  allora $A(\varphi)$ vale per ogni $\varphi \in \PROP$.
}

\pf{Dimostrazione}{
  Sia $X = \{\varphi \in \PROP \mid A(\varphi)\}$. Per (i)--(iii), $X$ soddisfa
  le tre clausole della definizione di $\PROP$. Poiché $\PROP$ è il più piccolo
  insieme con queste proprietà, $\PROP \subseteq X$, cioè $A$ vale per tutte le
  formule.
}

\subsection{Semantica}

Usiamo $1$ per ``vero'' e $0$ per ``falso''.

\dfn{Valutazione}{
  Una funzione $v : \PROP \to \{0,1\}$ è una \textbf{valutazione} se
  \begin{align*}
    v(\varphi \wedge \psi) &= \min(v(\varphi), v(\psi)), \\
    v(\varphi \vee \psi) &= \max(v(\varphi), v(\psi)), \\
    v(\varphi \to \psi) = 0 &\iff v(\varphi) = 1 \text{ e } v(\psi) = 0, \\
    v(\varphi \leftrightarrow \psi) = 1 &\iff v(\varphi) = v(\psi), \\
    v(\neg\varphi) &= 1 - v(\varphi), \\
    v(\bot) &= 0.
  \end{align*}
}

Le stesse regole come \textbf{tavole di verità}:
\begin{center}
\begin{tabular}{cc|ccccc}
  $\varphi$ & $\psi$ & $\neg\varphi$ & $\varphi\wedge\psi$ & $\varphi\vee\psi$ & $\varphi\to\psi$ & $\varphi\leftrightarrow\psi$ \\ \hline
  0 & 0 & 1 & 0 & 0 & 1 & 1 \\
  0 & 1 & 1 & 0 & 1 & 1 & 0 \\
  1 & 0 & 0 & 0 & 1 & \textbf{0} & 0 \\
  1 & 1 & 0 & 1 & 1 & 1 & 1
\end{tabular}
\end{center}

\thm{Estensione unica}{
  Se $v$ è una funzione dagli atomi in $\{0,1\}$ con $v(\bot) = 0$, esiste
  un'\textbf{unica} valutazione $\sem{\cdot}_v$ tale che $\sem{\varphi}_v =
  v(\varphi)$ per ogni atomo $\varphi$.
}

\nt{
  Si definisce $\sem{\cdot}_v$ per ricorsione sulla struttura della formula,
  usando le clausole della definizione di valutazione; l'unicità si dimostra
  per induzione. Conseguenza pratica: \textbf{una valutazione è determinata
  dai valori sugli atomi}. Nel caso proposizionale un ``mondo'' è quindi
  semplicemente un'assegnazione di $0/1$ alle proposizioni atomiche.
}

\mlenma{Dipendenza dai soli atomi presenti}{
  Se $v(p_i) = v'(p_i)$ per tutti i $p_i$ che compaiono in $\varphi$, allora
  $\sem{\varphi}_v = \sem{\varphi}_{v'}$.
}

\pf{Dimostrazione}{
  Induzione su $\varphi$. Se $\varphi$ è un atomo $p_i$ è l'ipotesi; se è
  $\bot$ entrambe valgono $0$. Se $\varphi = \varphi_1 \square \varphi_2$, gli
  atomi di $\varphi_1$ e $\varphi_2$ sono atomi di $\varphi$, quindi per
  ipotesi induttiva $\sem{\varphi_j}_v = \sem{\varphi_j}_{v'}$ per $j = 1,2$, e
  il valore di $\varphi$ si calcola dagli stessi valori con la stessa
  regola. Analogo per $\neg$.
}

\cor{Decidibilità}{
  Se $\varphi$ contiene $n$ atomi distinti, per sapere se è vera in tutte le
  valutazioni basta controllarne $2^n$ (le righe della tavola di verità).
  Quindi la validità proposizionale è \textbf{decidibile} (ma con costo
  esponenziale nel caso peggiore).
}

\dfn{Tautologia e conseguenza semantica}{
  \begin{enumerate}[label=(\roman*)]
    \item $\varphi$ è una \textbf{tautologia}, e si scrive $\models \varphi$,
          se $\sem{\varphi}_v = 1$ per ogni valutazione $v$.
    \item Sia $\Gamma$ un insieme di formule. $\Gamma \models \varphi$ sse per
          ogni $v$: se $\sem{\psi}_v = 1$ per ogni $\psi \in \Gamma$, allora
          $\sem{\varphi}_v = 1$. Si dice che $\varphi$ è
          \textbf{conseguenza semantica} di $\Gamma$.
  \end{enumerate}
  Convenzione: $\varphi_1, \dots, \varphi_n \models \psi$ sta per
  $\{\varphi_1, \dots, \varphi_n\} \models \psi$.
}

\nt{
  È esattamente la definizione dei lucidi, con ``mondo + interpretazione''
  sostituito da ``valutazione''.
}

\ex{Tautologie e conseguenze (Van Dalen)}{
  \begin{itemize}
    \item $\models \varphi \to \varphi$, \quad $\models \neg\neg\varphi \to
          \varphi$, \quad $\models \varphi \vee \psi \leftrightarrow \psi \vee
          \varphi$;
    \item $\varphi, \psi \models \varphi \wedge \psi$;
    \item $\varphi, \varphi \to \psi \models \psi$ \hfill (\emph{modus ponens});
    \item $\varphi \to \psi, \neg\psi \models \neg\varphi$ \hfill (\emph{modus tollens}).
  \end{itemize}
  Verifica di modus ponens: le valutazioni con $\sem{\varphi} = 1$ e
  $\sem{\varphi \to \psi} = 1$ sono solo quelle con $\sem{\psi} = 1$ (riga
  $1,1$ della tavola).
}

\ex{Un controesempio: affermazione del conseguente}{
  $\varphi \to \psi, \psi \not\models \varphi$: con $v(\varphi) = 0$ e $v(\psi)
  = 1$ le premesse sono vere e la conclusione è falsa. (È l'inferenza
  dell'\emph{abduzione}: plausibile ma non corretta.)
}

\subsection{Legami tra $\models$, connettivi ed equivalenza}

\mprop{Equivalenza e bi-implicazione}{
  $\varphi \equiv \psi$ sse $\models \varphi \leftrightarrow \psi$.
}

\pf{Dimostrazione}{
  Per il lemma di caratterizzazione, $\varphi \equiv \psi$ sse
  $\sem{\varphi}_v = \sem{\psi}_v$ per ogni $v$, sse
  $\sem{\varphi \leftrightarrow \psi}_v = 1$ per ogni $v$ (definizione di
  $\leftrightarrow$), sse $\models \varphi \leftrightarrow \psi$.
}

\thm{Teorema di deduzione semantico}{
  $\Gamma, \varphi \models \psi$ \ sse \ $\Gamma \models \varphi \to \psi$.
}

\pf{Dimostrazione}{
  ($\Rightarrow$) Sia $v$ che rende vere tutte le formule di $\Gamma$. Se
  $\sem{\varphi}_v = 0$, allora $\sem{\varphi \to \psi}_v = 1$. Se
  $\sem{\varphi}_v = 1$, $v$ rende vero $\Gamma \cup \{\varphi\}$, quindi
  $\sem{\psi}_v = 1$ e $\sem{\varphi\to\psi}_v = 1$.

  ($\Leftarrow$) Sia $v$ che rende vero $\Gamma \cup \{\varphi\}$. Allora
  $\sem{\varphi \to \psi}_v = 1$ e $\sem{\varphi}_v = 1$, quindi (unico caso
  falso dell'implicazione escluso) $\sem{\psi}_v = 1$.
}

\cor{Conseguenza e validità}{
  $\varphi_1, \dots, \varphi_n \models \psi$ sse
  $\models (\varphi_1 \wedge \dots \wedge \varphi_n) \to \psi$.
}

\thm{Conseguenza e insoddisfacibilità}{
  $\Gamma \models \varphi$ \ sse \ $\Gamma \cup \{\neg\varphi\}$ non ha
  alcuna valutazione che rende vere tutte le sue formule
  (è \textbf{insoddisfacibile}).
}

\pf{Dimostrazione}{
  $\Gamma \not\models \varphi$ sse esiste $v$ con $\Gamma$ tutto vero e
  $\sem{\varphi}_v = 0$, sse esiste $v$ con $\Gamma$ tutto vero e
  $\sem{\neg\varphi}_v = 1$, sse $\Gamma \cup \{\neg\varphi\}$ è
  soddisfacibile. Negando entrambi i lati si ottiene la tesi.
}

\nt{
  Quest'ultimo teorema è l'idea su cui si basano la \textbf{risoluzione} e
  quindi Prolog: per dimostrare che un goal segue dal programma, si aggiunge
  la sua \emph{negazione} e si cerca di derivare la contraddizione (la
  clausola vuota). Lo vedremo nelle prossime lezioni.
}

\subsection{Alcune equivalenze notevoli}

\thm{Leggi algebriche (Van Dalen 1.3.1, 1.3.4)}{
  Sono tautologie:
  \begin{itemize}
    \item associatività e commutatività di $\wedge$ e $\vee$;
    \item distributività: $\varphi \vee (\psi \wedge \sigma) \leftrightarrow
          (\varphi \vee \psi) \wedge (\varphi \vee \sigma)$ e
          $\varphi \wedge (\psi \vee \sigma) \leftrightarrow (\varphi \wedge
          \psi) \vee (\varphi \wedge \sigma)$;
    \item De Morgan: $\neg(\varphi \vee \psi) \leftrightarrow \neg\varphi \wedge
          \neg\psi$ e $\neg(\varphi \wedge \psi) \leftrightarrow \neg\varphi
          \vee \neg\psi$;
    \item doppia negazione: $\neg\neg\varphi \leftrightarrow \varphi$;
    \item $(\varphi \to \psi) \leftrightarrow (\neg\varphi \vee \psi)$;
    \item $(\varphi \leftrightarrow \psi) \leftrightarrow (\varphi \to \psi)
          \wedge (\psi \to \varphi)$;
    \item $\neg\varphi \leftrightarrow (\varphi \to \bot)$.
  \end{itemize}
}

\pf{Dimostrazione (De Morgan, le altre sono analoghe)}{
  $\sem{\neg(\varphi\vee\psi)} = 1 \iff \sem{\varphi\vee\psi} = 0 \iff
  \sem{\varphi} = \sem{\psi} = 0 \iff \sem{\neg\varphi} = \sem{\neg\psi} = 1
  \iff \sem{\neg\varphi \wedge \neg\psi} = 1$, per ogni valutazione.
}

\ex{Esercizio svolto: contrapposizione}{
  $(\varphi \to \psi) \equiv (\neg\psi \to \neg\varphi)$. Infatti
  \[
    \neg\psi \to \neg\varphi \equiv \neg\neg\psi \vee \neg\varphi
    \equiv \psi \vee \neg\varphi \equiv \neg\varphi \vee \psi
    \equiv \varphi \to \psi .
  \]
  Ogni passo usa una delle leggi sopra; il passaggio da un'equivalenza
  all'altra è lecito perché $\equiv$ è transitiva (teorema precedente) e
  perché sostituire una sottoformula con una equivalente dà una formula
  equivalente (\emph{teorema di sostituzione}, Van Dalen 1.2.6).
}

\ex{Esercizio svolto: Curry rivisto}{
  Esiste una valutazione con $\sem{X \leftrightarrow (X \to Y)} = 1$ e
  $\sem{Y} = 0$? Con $\sem{Y} = 0$ si ha $\sem{X \to Y} = 1 - \sem{X}$, quindi
  servirebbe $\sem{X} = 1 - \sem{X}$: impossibile. Dunque
  $X \leftrightarrow (X \to Y) \models Y$: l'``assunzione 1'' del paradosso di
  Curry implica già $Y$ da sola, ed è soddisfacibile solo nei mondi dove $Y$ è
  vero. Nessun paradosso.
}
